% LaTeX companion of hw08.typ. Both formats are editable.
\section{Homework 8 --- submitted work}

\chapter{Part A}

The assigned book exercises are 5.1: 45; 5.2: 14, 26; 5.3: 36; and 5.4: 26, 32.

\section{5.1 Exercise 45}

For \(A = \begin{pmatrix}
3 & 5 & 11 \\
5 & 9 & 20 \\
11 & 20 & 49
\end{pmatrix}\) and \(V = \text{span}\left( {\overrightarrow{v_{2}},\overrightarrow{v_{3}}} \right)\), put \(\text{proj}_{V{(\overrightarrow{v_{1}})}} = c_{2}\overrightarrow{v_{2}} + c_{3}\overrightarrow{v_{3}}\). Orthogonality gives

\(9c_{2} + 20c_{3} = 5,\quad 20c_{2} + 49c_{3} = 11,\)

so \(c_{2} = \frac{25}{41}\), \(c_{3} = - \frac{1}{41}\), and

\(\text{proj}_{V{(\overrightarrow{v_{1}})}} = \frac{25}{41}\overrightarrow{v_{2}} - \frac{1}{41}\overrightarrow{v_{3}}.\)

\section{5.2 Exercise 14}

For \(\begin{pmatrix}
1 \\
7 \\
1 \\
7
\end{pmatrix},\begin{pmatrix}
0 \\
7 \\
2 \\
7
\end{pmatrix},\begin{pmatrix}
1 \\
8 \\
1 \\
6
\end{pmatrix}\), the submitted Gram--Schmidt calculation gives

\(u_{1} = \frac{1}{10}\begin{pmatrix}
1 \\
7 \\
1 \\
7
\end{pmatrix},\quad u_{2} = \frac{1}{\sqrt{2}}\begin{pmatrix}
{- 1} \\
0 \\
1 \\
0
\end{pmatrix},\quad u_{3} = \frac{1}{\sqrt{2}}\begin{pmatrix}
0 \\
1 \\
0 \\
{- 1}
\end{pmatrix}.\)

It states that \(\left( {u_{1},u_{2},u_{3}} \right)\) is the orthonormal basis.

\section{5.2 Exercise 26}

For \(m_{1} = \begin{pmatrix}
2 \\
3 \\
0 \\
6
\end{pmatrix}\) and \(m_{2} = \begin{pmatrix}
4 \\
4 \\
2 \\
13
\end{pmatrix}\),

\(u_{1} = \frac{1}{7}\begin{pmatrix}
2 \\
3 \\
0 \\
6
\end{pmatrix},\quad u_{2} = \frac{1}{3}\begin{pmatrix}
0 \\
{- 2} \\
2 \\
1
\end{pmatrix},\)

and

\(Q = \begin{pmatrix}
\frac{2}{7} & 0 \\
\frac{3}{7} & {- \frac{2}{3}} \\
0 & \frac{2}{3} \\
\frac{6}{7} & \frac{1}{3}
\end{pmatrix},\quad R = \begin{pmatrix}
7 & 14 \\
0 & 3
\end{pmatrix}.\)

\section{5.3 Exercise 36}

For \(\begin{pmatrix}
\frac{2}{3} & \frac{1}{\sqrt{2}} & a \\
\frac{2}{3} & {- \frac{1}{\sqrt{2}}} & b \\
\frac{1}{3} & 0 & c
\end{pmatrix}\), the roles must be orthonormal. The page solves

\(\frac{4}{9} - \frac{1}{2} + ab = 0,\quad\frac{2}{9} + ac = 0,\quad\frac{2}{9} + bc = 0,\)

giving \(a = b = \frac{\sqrt{2}}{6}\) and \(c = - 2\frac{\sqrt{2}}{3}\).

\section{5.4 Exercise 26}

For \(A = \begin{pmatrix}
1 & 2 & 3 \\
4 & 5 & 6 \\
7 & 8 & 9
\end{pmatrix}\), \(b = \begin{pmatrix}
1 \\
0 \\
0
\end{pmatrix}\), the normal equation is

\(\begin{pmatrix}
66 & 78 & 90 \\
78 & 93 & 108 \\
90 & 108 & 126
\end{pmatrix}x = \begin{pmatrix}
1 \\
2 \\
3
\end{pmatrix}.\)

Its reduced echelon form is

\(\begin{pmatrix}
1 & 0 & {- 1} & \frac{7}{6} \\
0 & 1 & 2 & 1 \\
0 & 0 & 0 & 0
\end{pmatrix},\)

so

\(x^{\ast} = \begin{pmatrix}
{- \frac{7}{6} + t} \\
{1 - 2t} \\
t
\end{pmatrix} = \begin{pmatrix}
{- \frac{7}{6}} \\
1 \\
0
\end{pmatrix} + t\begin{pmatrix}
1 \\
{- 2} \\
1
\end{pmatrix}.\)

\section{5.4 Exercise 32}

For \(\left( {0,27} \right),\left( {1,0} \right),\left( {2,0} \right),\left( {3,0} \right)\), let \(f(x) = c_{0} + c_{1}x + c_{2}x^{2}\), with

\(A = \begin{pmatrix}
1 & 0 & 0 \\
1 & 1 & 1 \\
1 & 2 & 4 \\
1 & 3 & 9
\end{pmatrix},\quad b = \begin{pmatrix}
27 \\
0 \\
0 \\
0
\end{pmatrix}.\)

The source gives

\(\begin{pmatrix}
4 & 6 & 14 \\
6 & 14 & 36 \\
14 & 36 & 98
\end{pmatrix}x = \begin{pmatrix}
27 \\
0 \\
0
\end{pmatrix},\)

\(x^{\ast} = \begin{pmatrix}
\frac{513}{20} \\
{- \frac{567}{20}} \\
\frac{21}{4}
\end{pmatrix},\)

and

\(f^{\ast}(x) = \frac{513}{20} - \frac{567}{20}x + \frac{21}{4}x^{2} = 25.65 - 28.35x + 6.25x^{2}.\)

\chapter{Part B}

\section{Problem 1}

For \(\pi(v) = \sum_{i = 1}^{d}\frac{v \cdot v_{i}}{v_{i} \cdot v_{i}}v_{i}\):

\subsection{(a)}

If \(v_{i} \cdot v_{j} = 0\) for \(i \neq j\), then \(\left( {\frac{v_{1}}{\left\| v_{1} \right\|},\ldots,\frac{v_{d}}{\left\| v_{d} \right\|}} \right)\) is an orthonormal basis. Writing \(u_{i} = \frac{v_{i}}{\left\| v_{i} \right\|}\),

\(\pi(v) = \sum_{i = 1}^{d}\frac{v \cdot v_{i}}{\left\| v_{i} \right\|^{2}}v_{i} = \sum_{i = 1}^{d{({v \cdot u_{i}})}}u_{i},\)

so \(\pi\) is the orthogonal projection onto \(W\).

\subsection{(b)}

If the basis is not perpendicular, some \(v_{i} \cdot v_{j} = a \neq 0\). While \(\text{proj}_{W{(v_{i})}} = v_{i}\),

\(\pi\left( v_{i} \right) = v_{i} + \ldots + \frac{a}{\left\| v_{j} \right\|^{2}}v_{j} + \ldots \neq v_{i}\)

by linear independence. Thus \(\pi\) is not \(\text{proj}_{W}\).

\section{Problem 2}

For the set \(O_{n}\) of orthogonal matrices:

\begin{itemize}
\tightlist
\item
  (a) False. \(I_{n} \in O_{n}\), but \(I_{n} + I_{n}\) is not orthogonal; the page gives \(\left( {I_{n} + I_{n}} \right)^{T{({I_{n} + I_{n}})}} = \begin{pmatrix}
  2 & 0 \\
  0 & 2
  \end{pmatrix}\), whereas its inverse is \(\begin{pmatrix}
  \frac{1}{2} & 0 \\
  0 & \frac{1}{2}
  \end{pmatrix}\).
\item
  (b) True. The composition of the orthogonal maps represented by \(A,B\) has standard matrix \(AB\), so \(AB \in O_{n}\).
\item
  (c) True by the same composition argument for \(A^{2}\).
\item
  (d) True. If \(A^{2}\) were orthogonal but \(A\) were not, with \(\left\| T_{A{(x)}} \right\| = a\left\| x \right\|\) and \(a \neq 1\), then \(\left\| {T_{A} \circ T_{A{(x)}}} \right\| = a^{2}\left\| x \right\| \neq \left\| x \right\|\), a contradiction.
\item
  (e) True. \(A \in O_{n}\) gives \(A^{T} = A^{- 1}\); \(A^{2} = I_{n}\) gives \(A = A^{- 1} = A^{T}\), so \(A\) is symmetric.
\end{itemize}

\section{Problem 3}

For an orthonormal basis \(\mathcal{B}\),

\(v = v_{1}b_{1} + \ldots + v_{r}b_{r},\quad w = w_{1}b_{1} + \ldots + w_{r}b_{r},\)

and the calculation on the page is

\(v \cdot w = \sum_{i = 1}^{r}v_{i}w_{i} = \lbrack v\rbrack_{\mathcal{B}} \cdot \lbrack w\rbrack_{\mathcal{B}}.\)

For two orthonormal bases \(\mathcal{B},\mathcal{C}\), \(S_{\mathcal{C}\rightarrow\mathcal{B}}\) has \(\left( {i,j} \right)\) entry \(c_{j} \cdot b_{i}\) and \(S_{\mathcal{B}\rightarrow\mathcal{C}}\) has transpose entry \(b_{i} \cdot c_{j}\). Hence

\(S_{\mathcal{C}\rightarrow\mathcal{B}} = S_{\mathcal{B}\rightarrow\mathcal{C}}^{T} = S_{\mathcal{B}\rightarrow\mathcal{C}}^{- 1},\)

so \(S_{\mathcal{B}\rightarrow\mathcal{C}}\) is orthogonal.

\section{Problem 4}

The submitted answers are:

\begin{itemize}
\tightlist
\item
  (a) True. \(\left( {\text{ker}\ A} \right)^{\perp} = \text{im}\ A^{T}\), from \(\text{ker}\left( A^{T} \right) = \left( {\text{im}\ A} \right)^{\perp}\) and double orthogonal complement.
\item
  (b) True. \(\text{ker}\ A = \text{ker}\left( {A^{T}A} \right)\), then rank-nullity gives \(\text{rank}\ A = \text{rank}\left( {A^{T}A} \right)\).
\item
  (c) True. \(\text{ker}\left( A^{T} \right) = \left( {\text{im}\ A} \right)^{\perp}\) and rank-nullity yield \(\text{rank}\ A = \text{rank}\left( A^{T} \right)\).
\item
  (d) True. With (b), (c), and \(\left( A^{T} \right)^{T} = A\), \(\text{rank}\left( {A^{T}A} \right) = \text{rank}\left( {AA^{T}} \right)\).
\item
  (e) False. \(\text{dim}\left( {\text{ker}\ A} \right) = m - \text{rank}\ A\) but \(\text{dim}\left( {\text{ker}\ AA^{T}} \right) = n - \text{rank}\ A\); if \(n \neq m\) they cannot be equal.
\end{itemize}

\section{Problem 5}

\subsection{(a)}

For \(X = \left\{ {x_{1},\ldots,x_{r}} \right\}\), \(Y = \left\{ {y_{1},\ldots,y_{s}} \right\}\) and \(X \perp Y\),

\(x = \sum a_{i}x_{i},\quad y = \sum b_{j}y_{j}\)

implies \(x \cdot y = 0\), because every \(x_{i} \cdot y_{j} = 0\). Thus \(\text{span}\ X \perp \text{span}\ Y\).

\subsection{(b)}

For a relation

\(\sum a_{i}x_{i} + \sum b_{j}y_{j} = 0,\)

if coefficients on both sides are nonzero then the equal nonzero vectors would belong to \(\text{span}\ X \cap \text{span}\ Y\). But (a) and WS 16 give this intersection as \(0\). Hence all coefficients vanish and \(X \cup Y\) is linearly independent.

\subsection{(c)}

A pairwise orthogonal set with fewer than \(n + 1\) members is not maximal: append \(0\), use Gram--Schmidt, and add a new unit vector. A set with more than \(n + 1\) members would make \(Y \cup \left\{ 0 \right\}\) linearly independent in \(\mathbb{R}^{n}\), which is impossible. Therefore a maximal pairwise orthogonal set has exactly \(n + 1\) elements.

\section{Problem 6}

\subsection{(a)}

The source applies Gram--Schmidt to full-rank \(A\) to obtain an orthonormal basis, reverses the chosen basis, and uses a QR factorization. It then sets

It takes the orthonormal basis matrix and the transpose of the upper factor, obtaining the stated QL factorization. The resulting lower factor is triangular with positive diagonal.

\subsection{(b)}

The finished submission stops after the construction for (a); no visible proof for part (b) appears in the source PDF.
